% !TEX encoding = UTF-8
% !TEX TS-program = pdflatex
% !TEX root = ../tesi.tex
% !TEX spellcheck = it-IT

%*************************************************************
\chapter{Introduzione}
\label{cap:introduzione}
%*************************************************************

\section{L'azienda}

\begin{figure}[htp]
\centering
\includegraphics[width=\textwidth/2]{../immagini/logo_zucchetti}
\caption{Logo di Zucchetti S.p.a.}
\end{figure}

La \myCompany è una software house con sede a Lodi, che si occupa di soluzioni complete per le aziende, professionisti(commercialisti, consulenti del lavoro, avvocati, curatori fallimentari, notai ecc.) e pubbliche amministrazioni(Comuni, Province, Regioni, Ministeri, società pubbliche ecc.).\\
Il gruppo Zucchetti è la prima azienda italiana in Europa con oltre 3300 addetti, 1100 partner e oltre 105000 clienti.\\
Le soluzioni principali proposte dall'azienda sono :\\
\begin{itemize}
\item Softwere: gestionali, per la sicurezza sul lavoro, analisi business ecc.
\item Hardware: per la rilevazione presenze, controllo accessi e controllo produzione.
\item Servizi: di outsourcing, cloud computing e data center.
\end{itemize}

\subsection{Portal Studio}
Lo stage si è svolto nella sede distaccata di Padova che si occupa di ricerca e sviluppo.
Tra i software che vengono sviluppati in questa sede è presente Portal Studio che consiste in una WEB application per la creazione di siti web.\\
L'applicazione offre all'utente un set completo di strumenti per la creazione di pagine web, permette la creazione e modifica in modo grafico della struttura HTML, la gestione degli stili tramite editor grafico per il CSS ed inoltre permette di gestire i dati provenienti da diversi tipi di database, il loro filtraggio e il binding con varie strutture HTML come liste e tabelle.\\
Il software risulta essere molto maturo e oltre alle funzionalità sopracitate permette anche la creazione di portlet e pagelet e altri elementi riutilizzabili e la gestione di risorse come dati in formato JSON.


\section{Il progetto}
Il progetto proposto dall'azienda per lo stage, nasce dal desiderio di aggiungere all'applicazione Portal Studio una nuova funzionalità che consiste nell'offrire all'utente la possibilità di inserire nelle proprie pagine HTML dei template già pronti e selezionabili da un insieme prestabilito.\\
Questo desiderio ha portato l'azienda ad interessarsi ai template engine come Mustache.js, HandlebarJS ecc.\\
Lo stage si divide in due parti.\\
La prima consisteva nello studio dei template e degli aspetti ad essi correlati, la seconda nella realizzazione di un editor che ne permettesse la visualizzazione e la modifica.\\
Le parti più rilevanti del lavoro svolto sono state sicuramente:
\begin{itemize}
	\item la scelta della libreria da utilizzare nello sviluppo del progetto, che è stata individuata analizzando le caratteristiche di varie librerie per il templateing, cercando fra queste quella che in miglior modo si adeguasse ai requisiti del progetto.
	Questa libreria è stata scelta sapendo che l'azienda l'avrebbe utilizzata per il proprio applicativo quindi i risultati dell'analisi effettuata sono stati discussi con il tutor aziendale;
	
	\item l'integrazione di plug-in JQuery all'interno del template, che ha richiesto la valutazione di varie soluzioni che permettessero di ottenere il risultato desiderato;
	
	\item la realizzazione della parte dell'editor relativa alla modifica dei dati, che permette di adattarsi a qualsiasi tipo di template HTML ed SVG e inoltre, nonostante i pochi tipi primitivi forniti da JavaScript, di individuare a quale elemento (immagine, URL, colore, ecc.) si riferiscano i dati dei template.
	  
\end{itemize}
\subsection{Prima parte}
La prima parte del progetto inizia con la realizzazione di qualche template prototipo, utile sia per studiare le possibilità della libreria scelta sia per avere un insieme di template da inserire nell'editor che è stato realizzato in seguito.\\
Durante questa parte del progetto l'attenzione è stata rivolta alla possibilità di realizzare template statici, dinamici, template come composizione di altri template (es. lista di contatti) e template che utilizzano SVG.\\
In seguito alla realizzazione dei template prototipo è stato effettuato uno studio sulla possibilità di rendere i template responsive cioè permetterne la visualizzazione sia su dispositivi desktop che mobile.\\
La prima parte si è conclusa con uno studio sulla possibilità di realizzare template che contenessero al loro interno plug-in JQuery e sulla gestione del caricamento delle librerie necessarie al funzionamento dei template all'interno della pagina HTML.

\subsection{Seconda parte}
La seconda parte del progetto consisteva nella realizzazione di un editor grafico che permette all'utente la selezione di un template da una lista prestabilita.
In seguito alla selezione del template desiderato quest'ultimo verrà visualizzato in un box dedicato e tramite un editor, che viene costruito sulla base dei dati editabili del template (i dati sono contenuti in un oggetto JSON che fa parte del template), è possibile vedere il comportamento del template durante la modifica dei suoi dati.	\\
Per determinati template, come quelli considerati composti, l'editor deve dare la possibilità di visualizzare direttamente l'oggetto JSON contenente i dati e permetterne la modifica.\\
Non essendo presente una struttura di beck-end proposta dall'applicazione Portal Studio, perché ancora in fase di valutazione, l'editor è stato sviluppato separatamente, l'insieme dei template consiste in un gerarchia di directory contenenti i vari elementi che compongono i template (file HTML, JSON, immagini e librerie) suddivise per categoria.\\
Il caricamento delle risorse viene eseguito tramite chiamate http-get request, non essendo presenti delle API fornite dall'azienda.

